\documentclass[11pt]{article}
\usepackage{amsmath,amssymb,amsthm,booktabs}
\usepackage[margin=1in]{geometry}
\newtheorem{theorem}{Theorem}
\newtheorem{lemma}[theorem]{Lemma}
\newtheorem{definition}[theorem]{Definition}
\title{Problem 872: Recursive Bracelets}
\author{Project Euler Solutions}
\date{}
\begin{document}
\maketitle

\section{Problem Statement}

This problem concerns counting equivalence classes of bead colorings on a cycle under the dihedral group (bracelets), with a recursive compositional structure. The necklace counting formula is:
\[
N(n,k) = \frac{1}{n}\sum_{d \mid n}\varphi(d)\,k^{n/d}.
\]

\section{Mathematical Foundation}

\begin{theorem}[Burnside's Lemma]
Let a finite group $G$ act on a finite set $X$. The number of orbits is
\[
|X/G| = \frac{1}{|G|}\sum_{g \in G} |X^g|.
\]
\end{theorem}

\begin{proof}
Count pairs $(g,x)$ with $g \cdot x = x$ in two ways. Over $G$: $\sum_{g}|X^g|$. Over $X$: $\sum_{x}|G_x|$. By orbit-stabilizer, $|G_x| = |G|/|G \cdot x|$, so
\[
\sum_{x \in X} |G_x| = |G| \sum_{\text{orbits}} 1 = |G| \cdot |\text{orbits}|. \qedhere
\]
\end{proof}

\begin{theorem}[Necklace Count]
The number of distinct necklaces with $n$ beads and $k$ colors is
\[
N(n,k) = \frac{1}{n}\sum_{d \mid n}\varphi(d)\,k^{n/d}.
\]
\end{theorem}

\begin{proof}
Apply Burnside to $C_n$ acting on $k$-colorings. Rotation by $j$ fixes $k^{\gcd(n,j)}$ colorings. Grouping by $d = n/\gcd(n,j)$: exactly $\varphi(d)$ values of $j$ satisfy $\gcd(n,j)=n/d$, yielding the formula.
\end{proof}

\begin{theorem}[Bracelet Count]
Under the dihedral group $D_n$:
\[
B(n,k) = \begin{cases}
\frac{1}{2}N(n,k) + \frac{1}{4}(k+1)k^{n/2} & n \text{ even},\\[4pt]
\frac{1}{2}N(n,k) + \frac{1}{2}k^{(n+1)/2} & n \text{ odd}.
\end{cases}
\]
\end{theorem}

\begin{proof}
$D_n$ has $2n$ elements: $n$ rotations contribute $n \cdot N(n,k)$ total fixed points. For $n$ odd, each of the $n$ reflections fixes $k^{(n+1)/2}$ colorings. For $n$ even, $n/2$ reflections through vertices fix $k^{n/2+1}$ each, and $n/2$ through edges fix $k^{n/2}$ each. Dividing by $2n$ gives the result.
\end{proof}

\begin{lemma}[Aperiodic Necklaces]
The count of aperiodic necklaces of length $n$ over $k$ symbols is
\[
L(n,k) = \frac{1}{n}\sum_{d \mid n}\mu(d)\,k^{n/d}.
\]
Furthermore, $N(n,k) = \sum_{d \mid n} L(d,k)$.
\end{lemma}

\begin{proof}
Every necklace of length $n$ has a unique primitive period $d \mid n$, giving $N(n,k) = \sum_{d \mid n} L(d,k)$. M\"obius inversion on the divisor lattice yields the formula for $L(n,k)$.
\end{proof}

\section{Algorithm}

Compute $N(n,k)$ by enumerating divisors of $n$ with Euler's totient and modular exponentiation. Apply the recursive bracelet structure as specified by the problem. Aggregate results modulo $M$ if required.

\section{Complexity Analysis}

\begin{itemize}
\item \textbf{Time:} $O(d(n) \cdot \log n)$ per necklace count, where $d(n)$ = number of divisors.
\item \textbf{Space:} $O(1)$ auxiliary per count; $O(n)$ if memoization is needed for recursion.
\end{itemize}

\section{Answer}

\[
\boxed{2903144925319290239}
\]

\end{document}
